% NEW 2026-09-18 (§5.2 readability pass; handoff 2026-09-17_ch5_s52_setup_critique.md
% §V3, §W). REPLACES tab_5-2a_factors_varied and tab_5-2b_factors_fixed, which are
% deleted: two tables became one because this repo's own evaluation conventions §c
% already ruled the form --- "every declared parameter is listed; each row says
% either the band it was swept over and what moved, or that it was held and why.
% ONE REGISTER, BOTH KINDS OF ROW." Reti's paired tables read as a pair only
% because their columns are identical, which is a workaround for what a single
% table does natively; Kim's Table 4 (one parameter table, a Type row-group) is the
% closer precedent.
% THE JUSTIFICATION COLUMN IS GONE, on Marc's ruling (2026-09-18): "I would just
% lose the column on each of them --- what do I need them for?" The motivation is
% now one high-level clause per prose unit (why the attacker is varied, why the
% defence is), which answers the question a reader actually asks; a per-row reason
% answered one nobody asked and cost most of a page.
% ROW GROUPS ARE THE SIMULATION'S MOVING PARTS (Marc's framing: "what are the big
% abstract moving parts"), and they are the same four objects the prose units
% declare, in the same order. Each element sits in ONE value column: the varied
% rows read straight down a single column, so the experiment is visible as a shape
% rather than described as one. An empty cell is the genre's own mark for absent
% (figure_table_conventions §e1).
% TERMINOLOGY: "deployment interval" applies the RATIFIED registry row (Marc,
% 2026-09-02: the defensive move as a scheduling event is a DEPLOYMENT; "mutation"
% survives only inside Zhang's cited phrasing) --- ch5 had drafted against it.
% "interval distribution" (was "timing regime") and "run length" (was "horizon")
% are PROPOSED rows awaiting Marc's ruling; see terminology.md. The three §5.4/§5.5
% figure captions still say "mutation interval" and are NOT changed here --- that
% is existing prose and one ruling away.
% HAND-SET; the generator pass (C39) is still owed and now has one table fewer to
% emit. Values transcribed from data/results/ch5_defended/run_corpus.py (PROFILES,
% CONDITIONS, INTERVALS, HORIZON), GEOMETRY and the defaults in
% src/mtdsim/l3_simulation/movement/run.py, MTD_DURATION and
% ATTACK_DURATION['PENALTY'] in mtdnetwork/data/constants.py, and the target set
% from movement/targeting.py. 60 000 s is Marc's placeholder (C36).
% Geometry tested against the real class and geometry: no overfull box; the
% tabular measures 453.3 pt against a 455.2 pt text width, body height 427 pt.
% [H] NOT [htbp]: with §5.3 still a placeholder there is no following prose for a
% float to settle against, and at [htbp] this table drifted past the §5.3 heading
% on the built page --- the same fault Table 5.1 was pinned for. REVERT to [htbp]
% when §5.3's prose lands (the preamble's standing note on ch5 placeholder floats).
% Ratify on read.
% RESTRUCTURED 2026-09-20 on Marc's rulings (handoff §AJ, §AJ8): ONE VALUE COLUMN,
% Parameter | Value, Brown's TABLE I form with Kim's row groups. The test for a
% row: a setting a replicator must make, in a word chapters 2 and 4 gave the
% reader. Gone: Geometry (a code constant; split into chapter 2's four nouns),
% Topology (folded into Seeds), Declared inputs (to the footnote), Sink retrace,
% Cost model/memory, Confusion penalty (model facts, ch4), Adaptive selector,
% Deployments per run. Durations STAY, per mechanism (Marc: that a deployment
% takes time is needed here). Varied rows first in each group; SEMICOLONS
% separate levels, COMMAS sit inside one held value. Caption's design sentence
% now matches run_corpus.py build_jobs (objective is fully crossed). Values are
% malleable while the experiments solidify; 60 000 s left out until it is built.
% Everything above this note is the 2026-09-18 record, kept for the trail.
\begin{table}[htbp]
  \centering
  \caption[The experiment]{The settings of the experiment. A parameter with several values, separated by semicolons, is varied; the rest are held. Arm, objective, condition and deployment interval are run in every combination; the interval distribution is varied at the 200\,s interval under the targeted objective.}
  \label{tab:experiment}
  \tablestyle\setlength{\tabcolsep}{4pt}
  \begin{tabular}{@{}cP{3.4cm}P{10.9cm}@{}}
    \toprule
    & Parameter & Value \\
    \midrule
    & Hosts & 50 \\
    & Levels of depth & 4 \\
    & Subnets & 8 \\
    & Exposed endpoints & 5 \\
    \multirow{-5}{*}[-0.30cm]{\smash{\rotatebox{90}{\emph{Network}}}} & Target & the two database hosts, at the deepest level \\
    \midrule
    & Arm & the baseline attacker; the movement attacker on each of the four attack profiles; the movement attacker on the aggregate \\
    \multirow{-2}{*}[-0.46cm]{\smash{\rotatebox{90}{\emph{Attacker}}}} & Objective & targeted; opportunistic \\
    \midrule
    & Condition & no defence; each of the seven defence mechanisms alone; the random execution scheme; the alternative execution scheme \\
    & Deployment interval & 200\,s; 2\,000\,s \\
    & Interval distribution & quasi-periodic; exponential, same mean \\
    \multirow{-4}{*}[-0.55cm]{\smash{\rotatebox{90}{\emph{Defence}}}} & Deployment duration\textsuperscript{\dag} & complete topology shuffle 110\,s, host topology shuffle 100\,s, IP shuffle 100\,s, OS diversity 80\,s, port shuffle 70\,s, service diversity 70\,s, user shuffle 20\,s \\
    \midrule
    & Run length & 15\,000\,s \\
    \multirow{-2}{*}[-0.17cm]{\smash{\rotatebox{90}{\emph{Runs}}}} & Seeds & 1\,000 per combination, the same set on every arm \\
    \bottomrule
    \addlinespace[2pt]
    \multicolumn{3}{@{}p{0.96\textwidth}@{}}{\scriptsize \textsuperscript{\dag}~inherited from the simulator \citep{zhang2023}. The attacker model's inputs are at the declared values of Table~\ref{tab:parameter-register}. Versions: the partial tactic-to-verb mapping (Appendix~\ref{app:controller-mapping}); the failure-only weight set (Appendix~\ref{app:weight-sets}); ATT\&CK Enterprise v19.1.}\\
  \end{tabular}
\end{table}
